\scanpage{41}
\uline{Théorème} Soient $S$ un schéma. Pour toute forme $G$ de
$\mathrm{GP}(1)_S$ et tout \uline{mono}

\[
\varphi:\mathcal A_{5S}\lhook\joinrel\longrightarrow G
\]

% slanted
\margin{il suffit que ce soit \deleted{\unsure{non}} non trivial (donc mono) sur \unsure{tte fibre|les fibres}…}

soit

\[
\alpha(\varphi)=A(\varphi(\pi))
\]

où

\[
A:G\mathrel{\unsure{\rule{1.8em}{.4pt}|\longrightarrow}}\mathbb E^1_S
\]

est le ``\unsure{trace|hom.} invariant'' et où $\pi$ est la permutation
circulaire standart de $\mathcal A_5$. Ceci posé : le \deleted{cat. des}
foncteur \deleted{couples} $(G,\varphi)$ \deleted{\unsure{qui}} \inserted{$\longrightarrow\alpha(\varphi)$ de la catégorie-groupoïde des $(G,\varphi)$}
où $G$ est une forme de
$\mathrm{GP}(1)_S$ et $\varphi:\mathcal A_{5S}\hookrightarrow G$, \inserted{\deleted{avec les iso}}
\deleted{est \uline{discrète} (i.e. tout
autom. est une identité), et l'application
$(G,\varphi)\mapsto\alpha(\varphi)$ définit une bijection de l'ens de
ses classes d'iso, avec l'ens \deleted{des} $\mathrm{Hom}(S,M)$ des
$\alpha\in\Gamma(S,\mathcal O_S)$ satisfaisant
$\alpha^2+\alpha-1=0$}
dans la catégorie
discrète définie par l'ens des $\alpha\in\Gamma(S,\mathcal O_S)$ tels
que $\alpha^2+\alpha-1=0$,
est une équivalence de catégories. Ou encore :

\margin{On pose $\alpha'(\varphi)=-1-\alpha(\varphi)=-1/\alpha(\varphi)$}

La catégorie des $(G,\varphi)$ est \uline{rigide}, et pour
$S$ variable, le foncteur $S\longmapsto F(S)={}$ \deleted{ens} des
classes d'iso de couples $(G,\varphi)$ sur $S$, est représentable
grâce à $\varphi\mapsto\alpha(\varphi)$ par $M=\mathrm{Spec}\,\mathbb Z[T]/(T^2+T-1)$.

\uline{Cor} Soit $(G,\varphi)$ comme dessus, et soit $u$ un
autom. de $\mathcal A_5$. \deleted{Pour qu'on ait} \inserted{Si $u$ est intérieur on a}
$\alpha(\varphi\circ u)=\alpha(\varphi)$ i.e. $(G,\varphi)\sim(G,\varphi\circ u)$. Si $u$ n'est
intérieur sur aucune fibre, on a $\alpha(\varphi\circ u)=\alpha'(\varphi)$

% archive sheet number: 40
